\subsection{USB 相机采集}\label{sec:camera}

整条感知流水线的第一帧来自一只接在 USB3 端口上的相机。它属于 USB 视频类（UVC，一套把摄像头
的取景、格式协商与帧传输标准化的设备规范），因此无需专用驱动，任何符合该类的相机都以同一套
描述符与控制请求被识别和驱动。这只相机能交付两种格式，一种是 YUYV，即未经压缩、按 4:2:2 排布
的原始像素，取来只需做一次色彩空间转换；另一种是 MJPEG，每帧是一张独立的 JPEG，取来还要先解码。
两种格式后续都进入 RGA 与 NPU，区别只在前端是否多一步解码。

相机的像素流不像存储或网络那样按需搬运，它是连续不断、对时延敏感而对个别丢失不敏感的。USB 为这
类负载准备了等时传输（一种周期性、预先预约总线带宽、且不做重传的传输方式）。主机在配置端点时就为
它在每个服务周期里划定一份固定带宽，到点必送，送不下的那一份直接丢弃而非排队重发，于是时延有上界，
代价是允许偶发的帧内字节缺失。这正合视频流的需要，一帧迟到的画面对机器人毫无意义，宁可丢掉去取下
一帧新的。承载这条流的是 RK3588 上的 DWC3 控制器，它在主机模式下对软件呈现为一个标准的 xHCI（可
扩展主机控制器接口）主机，相机的枚举、端点配置与等时传输都经由 Starry 既有的 xHCI 驱动完成，本部
件因此复用这套主机栈，把工作集中在让等时端点在这块硬件上正确地周期取数。

\repfig{usb-path.png}{相机经 USB3 端口与 DWC3/xHCI 主机栈，以等时传输把帧送入“最新帧”槽，感知循环从槽里取走最新的一帧。}{fig:camera}

\paragraph{正确配置等时端点}
xHCI 用一份端点上下文描述每个端点应当如何被周期服务，其中的 \texttt{Interval} 字段决定控制器多久
为该端点取一次数据，它并不直接等于描述符里的 \texttt{bInterval}，而要按端点类型与连接速度换算。对
高速及更高速度的等时端点，正确的换算是 \texttt{Interval} 取 \texttt{bInterval} 减一，含义是每
$2^{\texttt{Interval}}$ 个 125\,$\mu$s 微帧服务一次。若图省事改用一个固定下界把它夹到 1，那么
\texttt{bInterval} 为 1 的高带宽端点会被编码成每两个微帧才服务一次，可用带宽当场减半，相机的内部
缓冲随即溢出，等时帧被成片截断。换算因此严格按规范分速度分类型计算，\texttt{usb-host} 的
\texttt{calculate\_xhci\_interval} 把这一点固定下来。同一段端点配置里，等时端点的错误计数被显式置零，
因为等时传输本就不重传，给它配置重试次数没有意义；中断与等时端点还按每微帧包数算出突发大小与一次
服务区间的最大载荷（max ESIT payload），一并写入上下文，控制器才会按相机申报的带宽真正预约总线。

\paragraph{用传输环逐包提交一帧}
一帧的数据不是一次提交，而是被切成多个等时数据包，每个包对应传输环（一段主机与设备共享、由 TRB
描述待办传输的环形缓冲）上的一个等时 TRB。\texttt{usb-host} 为每个包按相机的最大包大小与突发参数
算出突发计数与末段包数填入 TRB，逐包串成一个传输描述符提交给控制器，敲一次门铃即开始周期取数。
设备每完成一个包就回送一个传输事件，驱动据其完成码与残余长度算出该包实到字节，集齐一帧的所有包后
把整帧交还上层。这里有两处与等时语义直接相关的处理。其一，当周期环被取空，控制器会上报环下溢或环
上溢事件，它们并非某个传输的完成，而是“这一拍没数据可送”的提示，驱动按 Linux 的做法把这类事件当作
环空跳过，不让它们被误当成帧完成。其二，对实到字节短于申报长度的等时帧，驱动逐包归类为足额、偏短
或为零，并与缺失完成事件的包数分开统计，从而把“事件投递丢失”与“控制器侧限速或 FIFO 截断”两类成因
区分开，这是在板上诊断一条被截断的等时输入流时唯一能凭代码坐实的判据。

\paragraph{用户态只保留最新的一帧}
即便等时端点已正确取数，若取来的帧在用户态排队等待处理，时延仍会随队列累积而失控。采集侧因此不排队。
基于 libuvc 的采集线程独立运行，每收到一帧就进入一个互斥锁保护的“最新帧”槽，把帧序号、尺寸、格式与
像素数据整体写入该槽并覆盖其中原有的一帧。覆盖之前若上一帧尚未被取走，只把一个丢弃计数加一，绝不把
旧帧留下排队。感知循环从另一侧读取，先在锁内只看一眼槽里的帧序号，序号未变就直接返回、不做任何拷贝；
仅当序号变化才把这一帧拷出并打上采集时刻的时间戳。如此一来，无论感知与控制这一侧快慢，它每次拿到的
都是当下最新的一帧，旧帧被持续丢弃，端到端时延不被采集积压拖长。\textbf{这一“丢旧取新”的纪律由
\texttt{tennis-app} 的采集封装与 \texttt{Camera::poll} 共同维持，丢弃帧数也作为基准指标的一项被记录。}
